% ============================================================
%  Human & Fire Detection — Algorithm Evaluation Report
%  Smart Thermal System for Patient Safety Monitoring
%  Waveshare 26984 (80x62) Thermal Sensor Dataset
%  Afeka College of Engineering, Tel-Aviv
%
%  Compiled with: pdflatex (MiKTeX / TeX Live / Overleaf)
% ============================================================
\documentclass[11pt, a4paper]{article}

% ---- Packages -----------------------------------------------
\usepackage[a4paper, margin=2.5cm]{geometry}
\usepackage[T1]{fontenc}
\usepackage{lmodern}
\usepackage{booktabs}
\usepackage{multirow}
\usepackage{array}
\usepackage{amsmath}
\usepackage{amssymb}
\usepackage[table]{xcolor}
\usepackage{graphicx}
\usepackage{caption}
\usepackage{hyperref}
\usepackage{parskip}
\usepackage{enumitem}
\usepackage{float}

\providecommand{\arraybackslash}{\let\\\tabularnewline}
\makeatletter
\@ifundefined{do@row@strut}{\let\do@row@strut\relax}{}
\makeatother

% ---- Colours ------------------------------------------------
\definecolor{tpgreen}{RGB}{0,128,0}
\definecolor{tnblue}{RGB}{31,73,125}
\definecolor{fpred}{RGB}{192,0,0}
\definecolor{fnyellow}{RGB}{150,100,0}
\definecolor{headerblue}{RGB}{31,73,125}
\definecolor{cellgreenlight}{RGB}{198,239,206}
\definecolor{cellbluelight}{RGB}{221,235,247}
\definecolor{cellredlight}{RGB}{255,199,206}
\definecolor{cellyellowlight}{RGB}{255,235,156}

\hypersetup{colorlinks=true, linkcolor=headerblue, urlcolor=headerblue, citecolor=headerblue}

% ---- Confusion matrix macro ---------------------------------
% \confmat{posLabel}{negLabel}{TP}{FN}{FP}{TN}{posTot}{negTot}
\newcommand{\confmat}[8]{%
    \vspace{4pt}\begin{center}
    \renewcommand{\arraystretch}{1.8}%
    \begin{tabular}{lcc}
        \hline\rule{0pt}{2.4ex}%
        & \textbf{Pred:\ #1} & \textbf{Pred:\ #2} \\[2pt]\hline\rule{0pt}{2.6ex}%
        \textbf{Truth:\ #1} (#7\ fr.)
            & \cellcolor{cellgreenlight}\textcolor{tpgreen}{\textbf{TP}}\;$=$\;#3
            & \cellcolor{cellyellowlight}\textcolor{fnyellow}{\textbf{FN}}\;$=$\;#4 \\[4pt]
        \textbf{Truth:\ #2} (#8\ fr.)
            & \cellcolor{cellredlight}\textcolor{fpred}{\textbf{FP}}\;$=$\;#5
            & \cellcolor{cellbluelight}\textcolor{tnblue}{\textbf{TN}}\;$=$\;#6 \\[2pt]\hline
    \end{tabular}\end{center}\vspace{4pt}
}

% ---- Title --------------------------------------------------
\title{%
    {\large Afeka College of Engineering, Tel-Aviv}\\[0.4em]
    {\large Smart Thermal System for Patient Safety Monitoring}\\[0.8em]
    {\LARGE\bfseries Human \& Fire Detection — Algorithm Evaluation}\\[0.4em]
    {\normalsize Waveshare 26984 Thermal Sensor \quad 80$\times$62 px \quad
     17 Scenes \quad 8{,}259 Annotated Frame-Views}
}
\author{%
    Guy Chen \and Yaniv Blau \and Roy Lieberman\\[0.3em]
    \small Supervisor: Or Zilberberg \qquad Advisor: Dr.\@ Oshrit Hoffer
}
\date{June 2026}

% =============================================================
\begin{document}
\maketitle
\tableofcontents
\newpage

% =============================================================
\section{Overview}
% =============================================================

This report evaluates the human-detection and fire-detection algorithms of the
Smart Thermal System on the \textbf{new Waveshare 26984 dataset}: three thermal
sensors at 80$\times$62\,px and 8\,Hz, recorded across 17 scenes in the ward
setting. It supersedes the earlier MLX90640 (32$\times$24) evaluations and is
organised in two parts:

\begin{description}[leftmargin=2em, style=nextline]
    \item[Part I — Human Detection:] three detectors (AdaptiveThreshold,
        HOG+SVM, MobileNet-SSD) under a frame-level presence criterion and a
        strict IoU\,$\ge$\,0.5 localisation criterion.
    \item[Part II — Fire Detection:] two detectors (Otsu pipelined rule-based,
        and Fire-SVM) scored against warm-body hard negatives.
\end{description}

The headline result is that the higher-resolution, lower-noise Waveshare sensor
substantially improves every detector relative to the MLX90640 baseline. Fire
recall in particular more than doubles, and both fire detectors hold
\textbf{near-perfect precision} ($\le$\,1 false alarm across 6{,}886 fire-negative
frames that include all multi-person scenes).

\subsection{Sensor and dataset}

\begin{table}[H]
\centering
\caption{Sensor upgrade. The Waveshare provides $\approx$6.5$\times$ more pixels
and less than half the noise floor of the MLX90640.}
\renewcommand{\arraystretch}{1.3}
\begin{tabular}{lcc}
\toprule
 & \textbf{MLX90640 (old)} & \textbf{Waveshare 26984 (new)} \\
\midrule
Resolution     & 32$\times$24 (768 px) & 80$\times$62 (4{,}960 px) \\
Noise floor    & 1.5\,$^\circ$C        & 0.7\,$^\circ$C \\
Frame rate     & 8\,Hz                 & 8\,Hz \\
Annotated fire frames & 78             & 1{,}373 \\
Annotated person frames & $\sim$5{,}900 & 7{,}421 \\
\bottomrule
\end{tabular}
\end{table}

\subsection{Common methodology}

Every frame is reduced to a binary verdict and scored with a confusion matrix
and four derived metrics:
\begin{align*}
\text{Accuracy} &= \tfrac{TP+TN}{TP+TN+FP+FN}, &
\text{Precision} &= \tfrac{TP}{TP+FP}, \\
\text{Recall} &= \tfrac{TP}{TP+FN}, &
\text{F1} &= \tfrac{2\,\text{P}\cdot\text{R}}{\text{P}+\text{R}}.
\end{align*}
\textbf{Recall} is the primary safety metric (a missed fire or intruder is a
direct failure); \textbf{precision} guards against alert fatigue. All labels
follow the project convention: YOLO class 0 = fire / ignition source,
class 1 = person. Preprocessing (where used) is the Tateno pipeline
(Gaussian denoise $\to$ background subtraction $\to$ L1 rectification),
calibrated on even-indexed empty-room frames.

% #############################################################
\part*{Part I — Human Detection}
\addcontentsline{toc}{section}{Part I — Human Detection}
% #############################################################

\section{Setup (Human Detection)}

\subsection{Scenes and split}
The only truly empty recording is \texttt{empty\_room} (0 person boxes). Its
even-indexed frames calibrate one Tateno pipeline per channel; its odd-indexed
frames become true negatives. Every other labeled scene contributes
person-present positives. Within each (scene, channel) group the labeled frame
indices are sorted and split sequentially 80/10/10 into train/val/test, which
prevents temporal leakage while keeping every scene represented in every split.

\begin{table}[H]
\centering
\caption{Human-detection split statistics (frame-channel examples). Negatives are
empty frames: empty-room odd frames plus person-scene frames whose subject has
left the field of view.}
\renewcommand{\arraystretch}{1.3}
\begin{tabular}{lccc}
\toprule
\textbf{Split} & \textbf{Total} & \textbf{Person-present (Pos)} & \textbf{Empty (Neg)} \\
\midrule
Train      & 6{,}441 & 5{,}866 & 575 \\
Validation &    801  &    783  &  18 \\
Test       &    825  &    772  &  53 \\
\bottomrule
\end{tabular}
\end{table}

The dataset is strongly positive-skewed ($\approx$94\,\% person-present),
reflecting a ward in which people are present almost continuously; the small
empty-frame count makes the false-alarm rate a higher-variance statistic than
recall.

\subsection{Detectors}
\begin{enumerate}[leftmargin=2em]
    \item \textbf{AdaptiveThreshold} — rule-based local-contrast segmentation on
          the Tateno residual; no training. Resolution-parameterised.
    \item \textbf{HOG+SVM} — Histogram-of-Oriented-Gradients features with a
          linear SVM, trained on the train split. Resolution-parameterised.
    \item \textbf{MobileNet-SSD} — a deep single-shot detector
          (PyTorch, CUDA / RTX 2060 Super), 30 epochs. Resolution-fixed
          (a dedicated Waveshare checkpoint).
\end{enumerate}

\subsection{Success criteria}
\begin{description}[leftmargin=2.5em, style=nextline]
    \item[Criterion A — Presence:] a frame counts as a positive prediction if the
        detector emits \emph{any} box; for empty frames any box is a false alarm.
    \item[Criterion B — Localisation (IoU\,$\ge$\,0.5):] a positive frame is a TP
        only if a predicted box overlaps a ground-truth box with IoU\,$\ge$\,0.5.
\end{description}
The aggregate tables below report Criterion B (the strict criterion); the
\emph{mean IoU on positive frames} is reported alongside.

\section{Results (Human Detection)}

\subsection{Aggregate (test split)}

\begin{table}[H]
\centering
\caption{Human detection --- aggregate test-split metrics under the strict
localisation criterion (IoU\,$\ge$\,0.5). MobileNet-SSD dominates on every axis;
the classical detectors keep competitive recall but raise far more false alarms
on empty frames.}
\renewcommand{\arraystretch}{1.3}
\begin{tabular}{lcccccc}
\toprule
\textbf{Detector} & \textbf{Acc} & \textbf{Prec} & \textbf{Rec} & \textbf{F1}
    & \textbf{Mean IoU} & \textbf{False-alarm} \\
\midrule
AdaptiveThreshold & 76.6\% & 93.4\% & 80.7\% & 86.6\% & 56.5\% & 83.0\% \\
HOG+SVM           & 79.8\% & 92.7\% & 85.1\% & 88.7\% & 58.5\% & 98.1\% \\
\textbf{MobileNet-SSD} & \textbf{95.8\%} & \textbf{99.1\%} & \textbf{96.4\%}
    & \textbf{97.7\%} & \textbf{67.3\%} & \textbf{13.2\%} \\
\bottomrule
\end{tabular}
\end{table}

\begin{table}[H]
\centering
\caption{Human-detection confusion counts on the test split (772 person frames,
53 empty frames). \textcolor{tpgreen}{TP} = person detected,
\textcolor{fnyellow}{FN} = person missed, \textcolor{fpred}{FP} = false alarm on
empty, \textcolor{tnblue}{TN} = correct silence.}
\renewcommand{\arraystretch}{1.3}
\begin{tabular}{lcccc}
\toprule
\textbf{Detector}
    & \textcolor{tpgreen}{\textbf{TP}} & \textcolor{fnyellow}{\textbf{FN}}
    & \textcolor{fpred}{\textbf{FP}}   & \textcolor{tnblue}{\textbf{TN}} \\
\midrule
AdaptiveThreshold & 623 & 149 & 44 &  9 \\
HOG+SVM           & 657 & 115 & 52 &  1 \\
MobileNet-SSD     & 744 &  28 &  7 & 46 \\
\bottomrule
\end{tabular}
\end{table}

\subsection{Reading the results}

\begin{itemize}[leftmargin=2em]
    \item \textbf{MobileNet-SSD is the production choice.} It detects 744 of 772
        person frames (96.4\,\% recall), localises them well (67.3\,\% mean IoU),
        and false-alarms on only 7 of 53 empty frames. On the crowded scenes
        (\texttt{3ppl\_dance}, \texttt{3pp\_surprise}, \texttt{3pplhedroncolider},
        \texttt{2ppl\_fight}) it scores a perfect 100\,\% F1.
    \item \textbf{The classical detectors trade precision for simplicity.}
        AdaptiveThreshold (no training) and HOG+SVM both keep recall in the
        80--85\,\% range, but each emits a spurious box on nearly every
        \emph{empty} frame --- HOG fired on 52 of 53 (98\,\% false-alarm). With
        people present almost continuously this barely dents accuracy, but it
        makes them unsuitable as a standalone presence gate.
    \item \textbf{The shared hard case is \texttt{1\_man\_run}.} Its test tail is
        mostly the runner leaving the frame, so it is dominated by empty frames
        with residual motion heat; every detector struggles (MobileNet 40\,\%
        recall, classical detectors $<$\,22\,\% accuracy). This is the scene to
        target with more training data.
    \item \textbf{Resolution helps localisation.} Mean IoU on the 80$\times$62
        sensor reaches 67\,\% for MobileNet --- a box that is genuinely useful for
        the downstream contact / restricted-area logic, which the 32$\times$24
        MLX could not reliably provide.
\end{itemize}

% #############################################################
\part*{Part II — Fire Detection}
\addcontentsline{toc}{section}{Part II — Fire Detection}
% #############################################################

\section{Setup (Fire Detection)}

\subsection{Scenes and balance}
Fire-positive scenes are those containing $\ge$1 class-0 (fire) box:
\texttt{1person cig}, \texttt{heater\_in\_middle}, \texttt{man\_light\_cig}, and
\texttt{man\_light\_cig\_1}. All remaining 13 scenes are fire-negative and act as
\textbf{hard negatives}: they contain warm human bodies ($\sim$33\,$^\circ$C) and
a warm calibration room that the detector must not confuse with an ignition
source. Across all three channels the dataset is
\textbf{1{,}373 fire / 6{,}886 safe} frame-views (16.6\,\% positive).

\subsection{Detectors}
\begin{enumerate}[leftmargin=2em]
    \item \textbf{OtsuFireDetector} — the rule-based three-stage pipeline:
          absolute hot-pixel segmentation (hot-pixel bypass, $T_{\text{ign}}=45\,^\circ$C,
          $T_{\text{fire}}=60\,^\circ$C), a hierarchical area/temperature
          classifier, and a temporal mass-gradient confirmation stage. No
          training; evaluated on every frame.
    \item \textbf{FireSVMDetector} — Otsu blob features fed to an RBF SVM, trained
          on a 70/30 stratified split (per scene $\times$ camera $\times$
          fire/safe).
\end{enumerate}

\section{Results (Fire Detection)}

\subsection{Aggregate}
\begin{table}[H]
\centering
\caption{Fire detection — aggregate metrics. Precision is near-perfect for both
detectors; the Fire-SVM lifts recall by $\approx$19 points over the rule-based
pipeline at the cost of one false alarm.}
\renewcommand{\arraystretch}{1.3}
\begin{tabular}{lcccccc}
\toprule
\textbf{Detector} & \textbf{Acc} & \textbf{Prec} & \textbf{Rec} & \textbf{F1}
    & \textbf{False-alarm} & \textbf{Frames} \\
\midrule
Otsu (rule-based)        & 92.8\% & 99.9\% & 56.9\% & 72.5\% & 0.01\% & 8{,}259 (all) \\
Fire-SVM (RBF, test) & 95.9\% & 99.7\% & 75.5\% & 85.9\% & 0.05\% & 2{,}506 (30\%) \\
\bottomrule
\end{tabular}
\end{table}

\paragraph{Otsu confusion (all 8{,}259 frames).}
\confmat{Fire}{Safe}{781}{592}{1}{6{,}885}{1{,}373}{6{,}886}

\paragraph{Fire-SVM confusion (test split, 2{,}506 frames).}
\confmat{Fire}{Safe}{314}{102}{1}{2{,}089}{416}{2{,}090}

\subsection{Per-scene breakdown}
\begin{table}[H]
\centering
\caption{Fire detection, per fire-positive scene. Recall tracks how hot the boxed
object actually is: the lighter and heater scenes are detected reliably; the
cigarette scenes are recall-limited (see \S\ref{sec:fireceiling}).}
\renewcommand{\arraystretch}{1.3}
\begin{tabular}{llccccc}
\toprule
\textbf{Scene} & \textbf{Detector} & \textbf{Prec} & \textbf{Rec} & \textbf{F1}
    & \textbf{Mean IoU} & \textbf{Frames} \\
\midrule
\multirow{2}{*}{\texttt{man\_light\_cig\_1}} & Otsu & 100\% & 100\% & 100\% & 35.6\% & 285 \\
                                              & SVM  & 100\% & 100\% & 100\% & --    & 88 \\
\midrule
\multirow{2}{*}{\texttt{heater\_in\_middle}}  & Otsu & 100\% & 76.8\% & 86.9\% & 21.2\% & 1{,}137 \\
                                              & SVM  & 100\% & 92.9\% & 96.3\% & --    & 343 \\
\midrule
\multirow{2}{*}{\texttt{man\_light\_cig}}     & Otsu & 100\% & 32.6\% & 49.1\% & 3.9\% & 210 \\
                                              & SVM  & 100\% & 50.0\% & 66.7\% & --   & 65 \\
\midrule
\multirow{2}{*}{\texttt{1person cig}}         & Otsu & 100\% & 26.4\% & 41.7\% & 4.6\% & 1{,}029 \\
                                              & SVM  & 100\% & 50.0\% & 66.7\% & --   & 312 \\
\bottomrule
\end{tabular}
\end{table}

\subsection{The recall ceiling on cigarette scenes}
\label{sec:fireceiling}
The two cigarette scenes are recall-limited for a physical, not algorithmic,
reason. The annotator boxes the \emph{ignition-source object} (the cigarette) in
every frame it is visible, including the many frames where the tip is not glowing.
Measured peak temperatures inside the fire boxes confirm this:

\begin{table}[H]
\centering
\caption{Peak temperature inside fire boxes. The detector's $T_{\text{ign}}$ is
45\,$^\circ$C; a frame whose boxed object never exceeds it cannot be detected by
any absolute-threshold rule, which bounds achievable recall.}
\renewcommand{\arraystretch}{1.3}
\begin{tabular}{lcccc}
\toprule
\textbf{Scene} & \textbf{Median peak} & \textbf{$>$45\,$^\circ$C} & \textbf{$>$50\,$^\circ$C} & \textbf{$>$60\,$^\circ$C} \\
\midrule
\texttt{man\_light\_cig\_1} & 63.4\,$^\circ$C & 100\% & 92\% & 75\% \\
\texttt{heater\_in\_middle} & 54.5\,$^\circ$C & 77\%  & 62\% & 45\% \\
\texttt{man\_light\_cig}    & 42.7\,$^\circ$C & 33\%  & 19\% & 0\%  \\
\texttt{1person cig}        & 39.9\,$^\circ$C & 26\%  & 19\% & 10\% \\
\bottomrule
\end{tabular}
\end{table}

Recall closely tracks the ``$>$45\,$^\circ$C'' column: where the boxed object is
genuinely hot (\texttt{man\_light\_cig\_1}, \texttt{heater\_in\_middle}) recall is
high; where it is mostly a cold object (\texttt{1person cig}) recall is bounded
near the fraction of hot frames. This is the same label-vs-physics ceiling
identified on the MLX dataset --- but the better sensor pushes aggregate recall
from 9\% (Otsu) / 32\% (SVM) on MLX to \textbf{56.9\% / 75.5\%} here.

\section{Conclusions}

\begin{itemize}[leftmargin=2em]
    \item \textbf{Fire detection is precision-safe and recall-improved.} Both
        detectors raise essentially zero false alarms (1 FP total, on a warm
        calibration-room frame) while the Fire-SVM reaches 75.5\,\% recall ---
        more than double the MLX baseline. For sustained, genuinely hot sources
        (lighter, heater) recall is 93--100\,\%.
    \item \textbf{Residual misses are concentrated on cold cigarette frames},
        an annotation-vs-physics artifact rather than a detector failure. A
        cigarette whose tip is below 45\,$^\circ$C carries no thermal signature
        to detect.
    \item \textbf{Human and fire detection are complementary} in the runtime
        pipeline; the per-detector trade-offs in Part I inform which human
        detector to pair with the fire stage.
\end{itemize}

\end{document}
